\section{Off-library operating temperature benchmark}
\label{sec:off_library}

The measurements in Sections~\ref{sec:pwr} and
\ref{sec:godiva} take the fuel and moderator at \emph{library}
temperatures --- every nuclide's target $T$ lands exactly on
a column of the ENDF/B-VII.1 HDF5 file. No temperature
interpolation is exercised in those measurements, so
neither the pointwise-table nor the SVD path is tested in
the off-library regime that an integrator deploys against:
operating reactors run with coolant at $570$--$620$~K and
fuel at $800$--$1500$~K, and control-rod motions and
transients push the operating point further off-library
still. Continuous-temperature Monte Carlo codes
pseudo-interpolate between the two library columns
bracketing the target $T$: at each collision a random
$\xi$ is drawn and the lower-$T$ cross section is used
with probability
$(T_\text{hi} - T)/(T_\text{hi} - T_\text{lo})$, otherwise
the upper. Over a particle's lifetime this recovers the
exact linear-in-$T$ interpolant in
expectation~\cite{romano2015openmc}. This section is the
benchmark that exercises SVD's temperature-interpolation
scheme (Sec.~\ref{sec:ducru}).

This section runs a three-way comparison --- SVD, pointwise
table with OpenMC-style stochastic $T$-pick, and the
industry-standard ACE+WMP hybrid~\cite{josey2015wmp} --- on
the PWR pin cell of Sec.~\ref{sec:pwr} with a
\emph{uniform $+150$~K} offset applied to every nuclide's
library temperature: fuel runs at $1050$~K (between $900$
and $1200$~K library), moderator and clad at $750$~K
(between $600$ and $900$~K). Statistics are $10$ seeds
$\times 50\,000$ particles $\times 120$ active batches per
provider ($6 \cdot 10^7$ active histories per provider).

\subsection{Channel-consistent stochastic \texorpdfstring{$T$}{T}-pick}
\label{sec:off_library_bug}

A single collision's MicroXs build touches six channels
(elastic, inelastic, $n,2n$, $n,3n$, fission, capture) plus
a total for the pointwise table. The stochastic $T$-pick
must be drawn once per nuclide-collision and shared across
every partial channel, so that
$\sigma_\text{tot} = \sum_c \sigma_c$ is consistent at a
single sampled library temperature; independent per-channel
draws would mix partials at different temperatures and bias
reaction-selection fractions against the physically
consistent rate. The engine implements this by wrapping
each reaction's pointwise table in a
\texttt{StochTempTable} that exposes
\texttt{lookup\_at\_idx\_with\_pick(energy, idx, use\_hi)};
\texttt{TableXsProvider::lookup} draws $\xi$ once per
nuclide-collision via a thread-local splitmix64 and passes
the resulting \texttt{use\_hi} to every channel lookup. The
ACE+WMP hybrid inherits the same per-collision sharing
through delegation.

The diagnostic for channel-consistency is straightforward:
$k_\infty$ as a function of moderator temperature must be
monotonic. Measured on PWR with channel-consistent picking:

\begin{center}
\begin{tabular}{@{}lc@{}}
\toprule
Moderator offset & Table $k_\infty$ \\
\midrule
$+0$~K (on-library, $600$~K)     & $1.327$ \\
$+150$~K (off-library, $750$~K)  & $1.324$ \\
$+300$~K (on-library, $900$~K)   & $1.322$ \\
\bottomrule
\end{tabular}
\end{center}

\noindent monotonic in $T$ and within $\sim\!500$~pcm of
the bracketing endpoints, as expected from the
$\sim\!500$~pcm spectrum hardening across the
$600 \to 900$~K range.

\subsection{Three-point unity Ducru on PWR}
\label{sec:off_library_3t}

Section~\ref{sec:ducru_puoi} explains why the raw Ducru
weights must be renormalized to $\sum_j w_j = 1$ before
being applied to $V^\top$ columns of a $\log\sigma$ SVD.
A two-point unity-normalized variant on the bracketing
$|\mathcal{S}|=2$ subset is the minimal scheme:
\begin{equation}
\tilde w_\text{lo} = \frac{w_\text{lo}}{w_\text{lo}+w_\text{hi}},
\quad
\tilde w_\text{hi} = \frac{w_\text{hi}}{w_\text{lo}+w_\text{hi}},
\end{equation}
with $w_\text{lo}$, $w_\text{hi}$ from
Eq.~\ref{eq:ducru}. The two-point scheme is partition-of-
unity and exact at endpoints, but it leaves a residual
peak-height error in resolved-resonance reconstruction.
Extending $\mathcal{S}$ from two to three library
temperatures (nearest-three by $|T_i - T|$) and applying
the same unity-normalization
(Eq.~\ref{eq:ducru_unity}) captures a quadratic correction
to the Faddeeva-kernel reconstruction. The shipped engine
uses the three-point variant.

\paragraph{U-$238$ MT$=102$ rank probe.}
An offline experiment
(\texttt{scripts/u238\_capture\_rank\_probe.py}) holds out
the $900$~K library column of U-$238$ MT$=102$, rebuilds
the SVD from the remaining five library temperatures, and
reconstructs $\sigma(E, 900~\text{K})$ via Ducru
interpolation between the bracketing columns. On the
$6.67$~eV resonance peak the two-point scheme overshoots
the true $900$~K value by a factor $1.011$ (band-averaged
L2 error $1.8\%$); the three-point scheme at rank $3$
recovers $0.994$ of truth (L2 error $1.05\%$), and at
rank $5$ recovers $0.990$ of truth (L2 error $0.84\%$).
The peak-height correction propagates directly to PWR
$k_\infty$.

\paragraph{PWR three-way comparison at $+150$~K.}
$10$ seeds $\times 50\,000$ particles~$\times~120$ batches:

\begin{center}
\begin{tabular}{@{}lccc@{}}
\toprule
Provider         & $k_\infty$ & $\sigma_\text{seed}$ & ns/p \\
\midrule
SVD ($3$-pt unity Ducru) & $1.32593$ & $0.00059$ & $27\,298$ \\
SVD ($2$-pt unity, ablation) & $1.30707$ & $0.00037$ & $15\,685$ \\
Pointwise table    & $1.32304$ & $0.00050$ & $26\,451$ \\
ACE+WMP            & $1.32806$ & $0.00045$ & $27\,279$ \\
\bottomrule
\end{tabular}
\end{center}

The shipped three-point SVD is $213$~pcm
($\approx\!3.6\sigma_\text{seed}$) below ACE+WMP. The
two-point ablation row is reported for comparison: it
sits at $2094$~pcm vs WMP, illustrating the resonance
peak-height contribution that the two-point reconstruction
leaves on the table and that the three-point scheme
recovers --- a $9.8\times$ reduction from a single
interpolation-subset change.

SVD memory stays at $152.5$~MB against $199.7$~MB for
ACE+WMP ($22\%$ less); speed is $27\,298$ vs $27\,279$~ns/p
(parity within $\sigma$). At an off-library PWR operating
temperature SVD with three-point unity Ducru is smaller
than ACE+WMP, matches it on speed, and agrees with it on
$k_\infty$ to $\approx\!3.6\sigma_\text{seed}$ --- this is
the deployable operating regime for the SVD provider on
this geometry.

Table~\ref{tab:pwr_offlib} collects the full sweep.

\begin{table*}[t]
\centering
\caption{PWR pin cell at $+150$~K global offset
(fuel $900\to1050$~K, mod $600\to750$~K, clad $600\to750$~K),
three-way comparison, ten-seed / $50\,000$-particle /
$120$-batch statistics. The two-point row is reported as
an ablation against the shipped three-point variant.}
\label{tab:pwr_offlib}
\begin{tabular}{@{}lccccc@{}}
\toprule
Provider & $k_\infty$ & $\sigma_\text{seed}$ & Gap vs WMP (pcm)
         & ns/p & Memory (MB) \\
\midrule
SVD, $2$-pt unity Ducru (ablation) & $1.30707$ & $0.00037$ &
  $-2094$ & $15\,685$ & $152.5$ \\
\textbf{SVD, $3$-pt unity Ducru} & $\mathbf{1.32593}$ &
  $\mathbf{0.00059}$ & $\mathbf{-213}$ &
  $\mathbf{27\,298}$ & $\mathbf{152.5}$ \\
Pointwise + stochastic $T$-pick & $1.32304$ & $0.00050$ &
  $-502$ & $26\,451$ & $201.3$ \\
ACE+WMP (reference)              & $1.32806$ & $0.00045$ &
  $+0$   & $27\,279$ & $199.7$ \\
\bottomrule
\end{tabular}
\end{table*}

\subsection{Discrete-level rank reduction}
\label{sec:discrete_rank}

Discrete inelastic levels (MT$=51$--$91$) are stored at a
separate rank from the smooth reactions because a level's
Lorentzian shape in laboratory energy is weakly
temperature-dependent: one basis vector suffices.
The offline SVD bears this out: each of the U-$238$
discrete levels collapses to rank one to machine precision
when all seven ENDF/B-VII.1 library columns are included.
On Godiva, dropping discrete-level rank from $5$ to $1$
removes $\approx\!70\%$ of the engine's SVD working set
($556 \to 164$~MB, $41$ levels $\times$ three U isotopes)
and produces no measurable $k_\text{eff}$ shift at
ten-seed statistics ($1.00322 \to 1.00358$,
$\Delta = 36$~pcm, $\sigma = 60$~pcm per seed). The same
reduction is the single largest contributor to the
$22\%$ memory win in Table~\ref{tab:pwr_offlib}.

This is a per-reaction-rank policy only. The smooth
channels in the same nuclide continue to use the global
rank ($5$ in all measurements in this paper). The engine
exposes the discrete-level rank as a separate
\texttt{-{}-discrete-rank} command-line flag; the GPU path
retains a uniform rank across all reactions because the
CUDA kernel's basis-stride convention assumes one. This is
a CPU optimization for now.

\subsection{Negative result: simplex-projected QP}
\label{sec:off_library_qp}

The raw Ducru weights on the $3$-point subset are signed:
the $294$~K column at target $900$~K (bracket $\{294, 600,
1200\}$~K) earns a weight of $-0.12$ from Eq.~\ref{eq:ducru}
before normalization. Negative weights are uncomfortable
for a reconstruction that is physically a probability
mixture of reference-temperature cross sections, and the
Ducru~$2017$ paper explicitly defines a quadratic
programming variant (\cite{ducru2017}~§4) that enforces
both $\sum_j w_j = 1$ and $w_j \geq 0$. A cheap
approximation to that QP is Euclidean projection of the
raw-Ducru weight vector onto the probability
simplex~\cite{duchi2008simplex}: clip negative entries to
zero and renormalize until $\sum w = 1$. We implemented
it. On U-$238$~MT$=102$ held out at $900$~K,
bracket $\{294, 600, 1200\}$~K, the raw weights
$(-0.12, 0.62, 0.50)$ project to $(0.00, 0.56, 0.44)$.
L2 reconstruction error on the $6.67$~eV resonance band:

\begin{center}
\begin{tabular}{@{}lccc@{}}
\toprule
Weighting (rank 5) & L2 error & Peak ratio & Sign \\
\midrule
$2$-pt unity Ducru & $2.1\%$  & $0.998$ & signed \\
$3$-pt unity Ducru & $0.84\%$ & $0.990$ & signed \\
$3$-pt simplex QP  & $5.0\%$  & $1.043$ & non-neg \\
\bottomrule
\end{tabular}
\end{center}

The simplex-projected QP is \emph{worse} than either signed
scheme. The reason is that the simplex projection is
Euclidean-optimal in \emph{weight} space; the Ducru raw
weights are Euclidean-optimal in the \emph{Doppler-kernel}
norm (see~\cite{ducru2017}~§3).

\paragraph{Kernel-Gram QP (Ducru 2017~§4).}
We then implemented the actual constrained optimization from
§4 of the same paper: minimize
$(\mathbf{w} - \mathbf{w}_\text{raw})^\top \mathbf{G}\,
(\mathbf{w} - \mathbf{w}_\text{raw})$ subject to
$\mathbf{1}^\top \mathbf{w} = 1$ and $\mathbf{w} \geq 0$,
where $\mathbf{G}$ is the empirical kernel-Gram matrix
$\mathbf{X}^\top \mathbf{X}$ formed from the log-$\sigma$
columns of the three bracketing library temperatures. Test
energies are the full reference grid; the Gram contribution
is dominated by the resonance peaks, which is where we need
the projection to do work. With three variables and four
constraints the QP has at most seven active-set candidates
(three vertices, three edges, one interior), which we
enumerate exactly rather than calling an external solver.
An optional rank-aware variant projects each library column
onto the first $k$ singular modes before building $\mathbf{G}$,
so the QP optimizes the reconstruction error at the same
truncation rank the engine will actually produce. Both
variants yielded near-identical weights on U-$238$ MT$=102$.

The rank dependence of the result is the interesting finding.
On the same held-out $900$~K test:

\begin{table*}[t]
\centering
\small
\begin{tabular}{@{}lccc@{}}
\toprule
 & rank 3 & rank 5 & rank 5 \\
Weighting & L2 / peak & L2 / peak & weights \\
\midrule
$3$-pt unity Ducru (signed) &
 $1.05\%$ / $0.994$ &
 $\mathbf{0.84\%}$ / $\mathbf{0.990}$ &
 $\{-0.12, 0.62, 0.50\}$ \\
$3$-pt simplex QP (non-neg) &
 $5.02\%$ / $1.045$ &
 $5.01\%$ / $1.043$ &
 $\{0.00, 0.56, 0.44\}$ \\
$3$-pt kernel-Gram QP (non-neg) &
 $\mathbf{0.60\%}$ / $\mathbf{1.000}$ &
 $2.06\%$ / $0.983$ &
 $\{0.00, 0.36, 0.64\}$ \\
\bottomrule
\end{tabular}
\end{table*}

The kernel-Gram QP is the \emph{best} scheme at rank~3 — the
$6.67$~eV peak reconstructs to within $0.01\%$ of truth --- but
is strictly worse than the signed three-point unity at
rank~5, the production rank for the PWR measurements. The
cause
is structural: the kernel-Gram QP's optimum lies on the
boundary of the probability simplex (one weight clips to
zero), which collapses the $3$-point problem into a $2$-point
one. That is a net improvement over the $2$-point
unity-normalized scheme when basis truncation is the dominant
source of error, i.e.\ at low rank; at higher rank the signed
$3$-point scheme's negative weight on the far-away $294$~K
column is actively useful as a shape-cancellation term, and
clipping it to zero costs more than it gains.

Non-negativity of the interpolation weights is not a
universal property. At rank~5 with all 9 nuclides, the
unity-normalized signed scheme has more degrees of freedom
to cancel residual basis error; at low rank, or on a
single nuclide where the $6.67$~eV peak dominates
$k_\infty$, the kernel-Gram QP is the better choice. A
rank-adaptive dispatch is possible but outside the scope
of this work; the signed three-point unity scheme is the
production default at rank~5.

\subsection{CI guard and reproducibility}

A \texttt{scripts/pwr\_verdict.py} runner wraps the
end-to-end comparison and exits non-zero if any of the
following thresholds are violated, so the script can be
called from CI to guard against regression:
$k_\infty$ SVD vs ACE+WMP $\leq 800$~pcm,
memory ratio SVD/WMP $\leq 1.05$,
speed ratio SVD/WMP $\geq 0.97$. The machine-readable
report for the $+150$~K, three-point unity Ducru,
rank-$5$, discrete-rank-$1$ configuration is at
\texttt{outputs/pwr\_verdict\_3temp.json} and reproduces
Table~\ref{tab:pwr_offlib}.

%% ============================================================
